% ============================================
% Assistant Professor Cover Letter
% Farmingdale State College
% ============================================

\documentclass[11pt,letterpaper]{letter}

\usepackage{graphicx}
\usepackage{geometry}
\usepackage{microtype}
\usepackage[hidelinks]{hyperref}

\geometry{left=25mm,right=25mm,top=25mm,bottom=25mm}

\signature{Decory Edwards}
\address{Department of Economics \\ Johns Hopkins University \\ Baltimore, MD 21218 \\ \texttt{dedwar65@jh.edu}}

\begin{document}

\begin{letter}{Search Committee \\
Department of Economics \\
Farmingdale State College \\
Farmingdale, NY}

\opening{Dear Members of the Search Committee,}

I am writing to apply for the tenure-track Assistant Professor of Economics position at Farmingdale State College. I am a Ph.D. candidate in Economics at Johns Hopkins University, where I work under the supervision of Professors Christopher D. Carroll, Nicholas W. Papageorge, and Stelios Fourakis, and I expect to complete my degree in 2026. My research and teaching interests are in macroeconomics, inequality, and applied quantitative methods.

My research agenda focuses on understanding wealth inequality through the lens of heterogeneous returns to wealth and household behavior. In my job market paper, I use microdata from the Survey of Consumer Finances to estimate persistent heterogeneity in portfolio returns and embed these estimates in a structural macroeconomic model. I show that allowing for return heterogeneity substantially improves the model’s ability to match observed wealth concentration and produces realistic implications for consumption behavior and policy analysis. This work connects micro-level financial behavior to aggregate outcomes in a way that is both empirically grounded and policy relevant.

In related work, I use data from the Health and Retirement Study to study how trust and beliefs shape household investment outcomes. I find evidence of a nonlinear relationship between trust and realized returns, highlighting a behavioral channel through which social attitudes influence economic inequality. Together, my research emphasizes careful empirical measurement, transparent quantitative methods, and applications that speak directly to questions relevant for students pursuing careers in business, analytics, and applied economics.

\textbf{Teaching.} My teaching experience centers on undergraduate macroeconomics and applied policy courses. I have led weekly discussion sections and held regular office hours for Principles of Macroeconomics, Intermediate Macroeconomic Theory, and an upper-level macroeconomic policy seminar. My teaching philosophy emphasizes clarity, structure, and applied intuition. I aim to help students build confidence in quantitative reasoning by connecting economic theory to real data and practical decision-making. I am prepared to teach introductory and intermediate economics, as well as quantitative and data-oriented courses that support programs in Applied Economics, Business Management, and Sport Management.

Farmingdale State College’s emphasis on applied learning and student engagement strongly resonates with my approach to economics education. The College’s focus on preparing students for careers in business, technology, and analytics aligns well with my use of empirical data, computational tools, and real-world examples in the classroom. I am particularly interested in contributing to courses that integrate economics with data analysis and financial decision-making, including areas such as behavioral economics and the digital economy. My research and teaching naturally complement programs that emphasize practical skills while maintaining analytical rigor. I would also welcome opportunities to involve students in applied research projects and to participate in service activities that strengthen civic engagement and connections with community partners.

I would be excited to contribute to Farmingdale State College’s mission as a public institution committed to inclusive excellence, student success, and applied scholarship. Thank you for your time and consideration. I would welcome the opportunity to discuss my candidacy further.

\closing{Sincerely,}

\end{letter}
\end{document}
